% Injection-path decomposition. Bars are drawn from paths.dat; the dashed rule
% marks the all-off control, so every bar is read against what the frozen
% backbone alone achieves.
\begin{figure}[t]
\centering
\begin{tikzpicture}
\begin{axis}[draft, height=44mm, width=0.95\columnwidth,
  ybar, bar width=7pt,
  xlabel={}, ylabel={구조 순응도},
  symbolic x coords={both,skip,attn,off},
  xtick=data,
  xticklabels={양쪽, skip만\\(추론 시), attn만\\(추론 시), 양쪽 끔},
  x tick label style={font=\scriptsize, align=center},
  ymin=0, ymax=0.78,
  legend pos=north east,
  enlarge x limits=0.16,
  every axis plot/.append style={line width=0.4pt}]
  \addplot[fill=black!72, draw=black, area legend] table[x=label, y=sc0] {paths.dat};
  \addlegendentry{$r{=}0$}
  \addplot[fill=black!22, draw=black, area legend] table[x=label, y=sc3] {paths.dat};
  \addlegendentry{$r{=}0.3$}
  \draw[densely dashed, black!55, line width=0.4pt]
    (axis cs:both,\pathAllOffZero) -- (axis cs:off,\pathAllOffZero);
  \addlegendimage{line legend, densely dashed, black!55, line width=0.4pt}
  \addlegendentry{양쪽 끔, $r{=}0$}
\end{axis}
\end{tikzpicture}
\caption{두 주입 경로를 모두 켠 경우, 각각 단독, 그리고 둘 다 끈 경우. skip 잔차가
구조 신호의 대부분을 나르고, 교차 어텐션 단독으로는 거의 회복하지 못하며, 둘 다 끄면
동결 백본으로 떨어진다(FID \pathFidAllOff{}, 양쪽 켰을 때는 \pathFidBoth{}, skip만
\pathFidSkipOnly{}, 어텐션만 \pathFidAttnOnly{}). 점선은 $r{=}0$에서 양쪽을 끈 대조군이다.
경로는 학습된 주 모델에서 추론 시에 끈 것이며, 따로 학습한 skip 단독 모델은
표~\ref{tab:skiponly}에 있다. $r{=}0$에서는 대조군 대비 단일 경로 이득의 합이 양쪽 경로
이득과 대략 같고, $r{=}0.3$에서는 양쪽 경로의 이득이 단일 경로 이득의 합을 넘는다.}
\label{fig:paths}
\end{figure}
